\subsection{ADJOINT REPRESENTATION}

Let \(G\) be a Lie group. Consider the law of analytic left operation

\[
(g, g ^ {\prime}) \mapsto g g ^ {\prime} g ^ {- 1} = (\operatorname{Int} g) g ^ {\prime}
\]

of G into G. This law of operation defines, by no. 3, a bilinear mapping of \(\mathcal{T}^{(\infty)}(\mathrm{G})\times \mathcal{T}^{(\infty)}(\mathrm{G})\) into \(\mathcal{T}^{(\infty)}(\mathrm{G})\) , which we shall denote by \(\top\) in this no. By Proposition 13 of no. 3,

\[
(t * t ^ {\prime}) \top t ^ {\prime \prime} = t \top (t ^ {\prime} \top t ^ {\prime \prime})\tag{22}
\]

for all \(t, t', t''\) in \(\mathcal{T}^{(\infty)}(\mathrm{G})\) . By Proposition 14 (i) of no. 3,

\[
\varepsilon_ {g} \top t = (\operatorname{Int} g) _ {*} t\tag{23}
\]

for all \(g \in G\) and \(t \in \mathcal{T}^{(\infty)}(G)\) . In particular, the mapping \(t \mapsto \varepsilon_{g} \top t\) of \(\mathcal{T}^{(\infty)}(G)\) into \(\mathcal{T}^{(\infty)}(G)\) is an automorphism of the bigebra \(\mathcal{T}^{(\infty)}(G)\) . Its restrictions to \(U(G)\) , \(U_{s}(G)\) , \(L(G)\) are denoted by \(\mathrm{Ad}_{\mathrm{U(G)}}(g)\) , \(\mathrm{Ad}_{\mathrm{U_s(G)}}(g)\) , \(\mathrm{Ad}_{\mathrm{L(G)}}(g)\) . We often write \(\mathrm{Ad}(g)\) instead of \(\mathrm{Ad}_{\mathrm{L(G)}}(g)\) when no confusion is possible. By (23), \(\mathrm{Ad}(g)\) is the tangent mapping at \(e\) to \(\operatorname{Int}(g)\) . It is an automorphism of the normable

Lie algebra \(\mathbf{L}(\mathbf{G})\) . When \(\mathbf{K}\) is of characteristic 0, \(\operatorname{Ad}_{\mathbf{U}(\mathbf{G})}(g)\) is the unique automorphism of \(\mathbf{U}(\mathbf{G})\) which extends \(\operatorname{Ad}(g)\) .

If \(\phi\) is a morphism of the Lie group G into a Lie group H, then

\[
\phi_ {*} (t \top t ^ {\prime}) = \phi_ {*} (t) \top \phi_ {*} (t ^ {\prime})\tag{24}
\]

for all \(t, t'\) in \(\mathcal{T}^{(\infty)}(\mathrm{G})\) ; this follows from Proposition 15 of no. 3.

PROPOSITION 43. Let \(t, u\) be in \(\mathcal{T}^{(\infty)}(\mathrm{G})\) . Let \(\sum_{i=1}^{n} t_i \otimes t'_i\) be the image of \(t\) under the coproduct. Then

\[
t \top u = \sum_ {i = 1} ^ {n} t _ {i} * u * t _ {i} ^ {\prime \vee}.
\]

By definition, \(t \top u\) is the image of \(t \otimes u\) under the mapping \((g, g') \mapsto gg'g^{-1}\) of \(G \times G\) into \(G\) . Now this mapping is obtained by composing the following mappings:

\[
\begin{array}{l l} \alpha \colon (g, g ^ {\prime}) \mapsto (g, g, g ^ {\prime}) & \text {of G \times G into G \times G \times G} \\ \beta \colon (g, g ^ {\prime}, g ^ {\prime \prime}) \mapsto (g, g ^ {\prime - 1}, g ^ {\prime \prime}) & \text {of G \times G \times G into G \times G \times G} \\ \gamma \colon (g, g ^ {\prime}, g ^ {\prime \prime}) \mapsto g g ^ {\prime \prime} g ^ {\prime} & \text {of G \times G \times G into G}. \end{array}
\]

On the other hand:

\[
\begin{array}{c} \alpha_ {*} (t \otimes u) = \sum_ {i = 1} ^ {n} (t _ {i} \otimes t _ {i} ^ {\prime}) \otimes u = \sum_ {i = 1} ^ {n} t _ {i} \otimes t _ {i} ^ {\prime} \otimes u \\ \beta_ {*} (\sum_ {i = 1} ^ {n} t _ {i} \otimes t _ {i} ^ {\prime} \otimes u) = \sum_ {i = 1} ^ {n} t _ {i} \otimes t _ {i} ^ {\prime \vee} \otimes u \\ \gamma_ {*} (\sum_ {i = 1} ^ {n} t _ {i} \otimes t _ {i} ^ {\prime \vee} \otimes u) = \sum_ {i = 1} ^ {n} t _ {i} * u * t _ {i} ^ {\prime \vee}. \end{array}
\]

COROLLARY 1. Let \(u \in \mathrm{L}(\mathbf{G})\) and \(u' \in \mathcal{T}^{(\infty)}(\mathbf{G})\) . Then \(u \top u' = u * u' - u' * u\) .

The image of u under the coproduct is \(u \otimes \varepsilon_{e} + \varepsilon_{e} \otimes u\) , whence

\[
u \top u ^ {\prime} = u * u ^ {\prime} * \varepsilon_ {e} + \varepsilon_ {e} * u ^ {\prime} * u ^ {\vee} = u * u ^ {\prime} - u ^ {\prime} * u.
\]

COROLLARY 2. Let \(t \in \mathcal{T}^{(\infty)}(\mathrm{G})\) and \(g \in \mathrm{G}\) . Then \(\varepsilon_g \top t = \varepsilon_g * t * \varepsilon_g^{-1}\) . If \(t \in \mathrm{L}(\mathrm{G})\) , then \(\varepsilon_g \top t = gtg^{-1}\) (where the latter product is evaluated in the group \(\mathrm{T}(\mathrm{G})\) ).

The image of \(\varepsilon_{g}\) under the coproduct is \(\varepsilon_{g} \otimes \varepsilon_{g}\) .

COROLLARY 3. Let \(a \in \mathrm{L}(\mathrm{G})\) . The vector field defined by \(a\) and the left operation \(g \mapsto \operatorname{Int} g\) of \(\mathrm{G}\) on \(\mathrm{G}\) is the field \(\mathbf{R}_a - \mathbf{L}_a\) .

The value of this field at \(g\) is

\[
\begin{array}{r l} a \top \varepsilon_ {g} = a * \varepsilon_ {g} - \varepsilon_ {g} * a & (\text { Corollary   1 }) \\ = (R _ {a}) _ {g} - (L _ {a}) _ {g} & (\text { Definition   5 }). \end{array}
\]

For all \(g \in G\) and all \(t \in \mathrm{L}(\mathrm{G})\) ,

\[
(\mathrm{Ad} g) (t) = \varepsilon_ {g} \top t = \varepsilon_ {g} * t * \varepsilon_ {g ^ {- 1}} = g t g ^ {- 1}.\tag{25}
\]

Since Ad \(g = \mathrm{T}_{e}(\mathrm{Int} g)\) , Proposition 42 of no. 11 proves that Ad is an analytic linear representation of G on the normable space \(\mathrm{L}(\mathrm{G})\) .

DEFINITION 8. The representation Ad of G on L(G) is called the adjoint representation of G.

PROPOSITION 44. For all \(a \in \mathbf{L}(\mathbf{G})\) ,

\[
(\mathrm{L} (\mathrm{Ad})) (a) = \mathrm{ad} _ {\mathrm{L} (\mathrm{G})} a.
\]

Let \(b \in \mathrm{L}(\mathrm{G})\) . By Proposition 42 (ii) of no. 11 and Corollary 3 to Proposition 43,

\[
(\mathrm{L} (\mathrm{Ad})) (a). b = - [ \mathrm{R} _ {a} - \mathrm{L} _ {a}, \mathrm{L} _ {b} ] (e).
\]

Now \(R_{a} \circ L_{b} = L_{b} \circ R_{a}\) (no. 6, Proposition 23 (ii)), whence \([R_{a}, L_{b}] = 0\) ; then, using Proposition 23 (ii),

\[
(\mathrm{L} (\mathrm{Ad})) (a). b = [ \mathrm{L} _ {a}, \mathrm{L} _ {b} ] (e) = \mathrm{L} _ {[ a, b ]} (e) = [ a, b ] = (\mathrm{ad} _ {\mathrm{L} (\mathrm{G})} a) b.
\]

PROPOSITION 45. Suppose that G is finite-dimensional and that K is of characteristic 0. Let s be an integer \(\geqslant0\) . Then the mapping \(\pi: g \mapsto \operatorname{Ad}_{\mathrm{U}_{s}(\mathbf{G})}(g)\) is an analytic linear representation of G on \(\mathrm{U}_{s}(\mathrm{G})\) and \(\mathrm{L}(\pi)a = \operatorname{ad}_{\mathrm{U}_{s}(\mathbf{G})}a\) for all \(a \in \mathrm{L}(\mathrm{G})\) .

The linear representation \(\pi\) is a quotient of \(\bigoplus_{r=0}^{\infty} \mathrm{T}^{r}(\mathrm{Ad})\) and is hence analytic. For \(a \in \mathrm{L}(\mathrm{G})\) and \(x_{1}, x_{2}, \ldots, x_{s}\) in \(\mathrm{L}(\mathrm{G})\) ,

\[
\begin{array}{l l} (\mathrm{L} (\pi) a) (x _ {1} x _ {2} \dots x _ {s}) = \sum_ {i = 1} ^ {s} x _ {1} \dots (\mathrm{L} (\mathrm{Ad}) a. x _ {i}) \dots x _ {s} & \text {(Proposition 41)} \\ = \sum_ {i = 1} ^ {s} x _ {1} \dots ([ a, x _ {i} ]) \dots x _ {s} & \text {(Proposition 44)} \\ = (\mathrm{Ad} _ {\mathrm{U} _ {s} (\mathrm{G})} a) (x _ {1} x _ {2} \dots x _ {s}). \end{array}
\]

PROPOSITION 46. Let \(h \in \mathbf{G}\) , \(x \in \mathrm{T}_h(\mathbf{G})\) and \(a \in \mathrm{L}(\mathbf{G})\) . Let \(\phi\) be the mapping \((g, g') \mapsto gg'g^{-1}\) of \(\mathbf{G} \times \mathbf{G}\) into \(\mathbf{G}\) . The image \(y\) of \((a, x) \in \mathrm{T}_e(\mathbf{G}) \times \mathrm{T}_h(\mathbf{G})\) under \(\mathrm{T}_{(e,h)}(\phi)\) is \(y = x + h((\mathrm{Ad}h^{-1})a - a)\) .

We have

\[
\begin{array}{r l} & y = (\mathrm{T} _ {(e, h)} \phi) (a \otimes \varepsilon_ {h} + \varepsilon_ {e} \otimes x) \\ & \quad = a \top \varepsilon_ {h} + \varepsilon_ {e} \top x \\ & \quad = a * \varepsilon_ {h} - \varepsilon_ {h} * a + x \\ & \quad = h ((\mathrm{Ad} h ^ {- 1}) a) - h a + x. \end{array}
\]

PROPOSITION 47. Let G be a Lie group, H and E Lie subgroups of G and suppose that \(hEh^{-1} = E\) for all \(h \in H\) . Then \(\mathcal{T}^{(\infty)}(H) \top \mathcal{T}^{(\infty)}(E) \subset \mathcal{T}^{(\infty)}(E)\) . In particular, \(\operatorname{Ad}(H)(L(E)) \subset L(E)\) and \([L(H), L(E)] \subset L(E)\) .

If \(t \in \mathcal{T}^{(\infty)}(\mathrm{H})\) and \(t' \in \mathcal{T}^{(\infty)}(\mathrm{E})\) , then \(t \otimes t' \in \mathcal{T}^{(\infty)}(\mathrm{H} \times \mathrm{E})\) and the image of \(\mathrm{H} \times \mathrm{E}\) under the mapping \((g, g') \mapsto gg'g^{-1}\) is contained in \(\mathrm{E}\) .

PROPOSITION 48. Let G be a Lie group and H and E Lie subgroups of G. Suppose that G is, as a Lie group, the semi-direct product of H by E. Let \(\rho\) be the linear representation \(g \mapsto (\mathrm{Ad} g) \mid \mathrm{L}(E)\) of the Lie group G on \(\mathrm{L}(E)\) (cf.~Proposition 47) and let \(\sigma\) be the restriction of \(\rho\) to H. Then:

\begin{enumerate}
\def\labelenumi{(\roman{enumi})}
\item
  \(\mathbf{L}(\mathbf{G})\) is the topological direct sum of \(\mathbf{L}(\mathbf{H})\) and \(\mathbf{L}(\mathbf{E})\) ;
\item
  L(H) is a subalgebra of L(G) and L(E) is an ideal of L(G);
\item
  \(\mathbf{L}(\sigma)\) is a linear representation of \(\mathbf{L}(\mathbf{H})\) on the Lie algebra of derivations of \(\mathbf{L}(\mathbf{E})\) ;
\item
  L(G) is the semi-direct product of L(H) by L(E) defined by L(σ) (Chapter I, § 1, no. 8).
\item
  is obvious and (ii) follows from Proposition 47. \(\mathbf{L}(\sigma) = \mathbf{L}(\rho) \mid \mathbf{L}(\mathbf{H})\) . Now by Propositions 40 (no. 11) and 44 (no. 12), \(\mathbf{L}(\rho)(t)\) is, for all \(t \in \mathbf{L}(\mathbf{G})\) , the restriction of \(\mathrm{ad}_{L(\mathbf{G})} t\) to \(\mathbf{L}(\mathbf{E})\) . This proves (iii). Using (i) and (ii), this also proves (iv).
\end{enumerate}

COROLLARY. Let G be a Lie group. Let \(\mathrm{T}_{e}(\mathrm{G})\) be given its unique commutative Lie algebra structure. Let \(\tau\) be an adjoint representation of \(\mathrm{L}(\mathrm{G})\) . Then the Lie algebra of \(\mathrm{T}(\mathrm{G})\) is the semi-direct product of \(\mathrm{L}(\mathrm{G})\) by \(\mathrm{T}_{e}(\mathrm{G})\) defined by \(\tau\) . In other words, for \(x, x'\) in \(\mathrm{L}(\mathrm{G})\) and \(y, y'\) in \(\mathrm{T}_{e}(\mathrm{G})\) ,

\[
[ (x, y), (x ^ {\prime}, y ^ {\prime}) ] = ([ x, x ^ {\prime} ], [ x, y ^ {\prime} ] + [ y, x ^ {\prime} ])
\]

(where the bracket on the left is evaluated in \(\mathbf{L}(\mathbf{T}(\mathbf{G}))\) and the brackets on the right in \(\mathbf{L}(\mathbf{G})\) ).

This follows from Proposition 48 and Proposition 6 of § 2, no. 2.

PROPOSITION 49. Let A be a complete normable unital associative algebra. We identify A with L(A*). Then, if \(g \in \mathbf{A}^*\) and \(y \in \mathbf{A}\) , (Ad g) \(y = gyg^{-1}\) .

Recall that \(\operatorname{Ad} g = \mathrm{T}_{1}(\operatorname{Int} g)\) . Let \(u_{g}\) be the mapping \(x \mapsto gxg^{-1}\) of A into A. The identity chart of \(\mathrm{A}^{*}\) into A transforms \(\operatorname{Int} g\) into \(u_{g} \mid \mathrm{A}^{*}\) . The tangent mapping at each point of \(\mathrm{A}^{*}\) to this mapping is equal to \(u_{g}\) , whence the proposition.

COROLLARY. For all \(g \in \mathbf{A}^*\) , let \(i(g)\) be the automorphism \(y \mapsto \mathrm{gyg}^{-1}\) of \(\mathbf{A}\) , so that \(i\) is an analytic linear representation of \(\mathbf{A}^*\) on \(\mathbf{A}\) . For all \(z \in \mathrm{L}(\mathbf{A}^*) = \mathbf{A}\) , \(\mathrm{L}(i)z\) is the inner derivation \(y \mapsto zy - yz\) of \(\mathbf{A}\) .

This follows from Propositions 49 and 44.

